%% main.tex — IP&M MMLowRes Special Issue submission
%% Working title: Tonal-Perceptual Caching for On-Device Vietnamese Speech Translation
%% Template: Elsevier cas-dc (single-column, double-spaced for initial submission)
%% Double-blind: all author identifying info is REDACTED. Restore before camera-ready.
%%
%% Build: pdflatex main && bibtex main && pdflatex main && pdflatex main
%% (see README.md in this directory)

\documentclass[a4paper,fleqn]{cas-dc}

\usepackage[authoryear,longnamesfirst]{natbib}

% Author macros
\def\tsc#1{\csdef{#1}{\textsc{\lowercase{#1}}\xspace}}
\tsc{WGM}
\tsc{QE}

% Theorem environments (uncomment if needed)
%\newtheorem{theorem}{Theorem}
%\newtheorem{lemma}[theorem]{Lemma}
%\newdefinition{rmk}{Remark}

\begin{document}
\let\WriteBookmarks\relax
\def\floatpagepagefraction{1}
\def\textpagefraction{.001}

\shorttitle{Tonal-Perceptual Caching for On-Device Vietnamese S2ST}
\shortauthors{Anonymous}

\title[mode=title]{Tonal-Perceptual Caching for On-Device Speech Translation of Vietnamese: A Typology-Aware Fingerprint with Two-Stream Voting}

\tnotemark[1]
\tnotetext[1]{Submitted to Information Processing \& Management, Special Issue ``Multilingual and Multimodal Low-Resource Language Processing'' (VSI: MMLowRes). Double-blind submission.}

%% ===== DOUBLE-BLIND: author block redacted. Replace before camera-ready. =====
\author[1]{Anonymous Author}
\cormark[1]
\fnmark[1]
\ead{anonymous@example.org}
\credit{Anonymized for double-blind review}

\affiliation[1]{organization={Anonymous Institution},
            addressline={},
            city={},
            postcode={},
            state={},
            country={}}

\cortext[1]{Corresponding author}
\fntext[1]{Author identifying details withheld for double-blind review.}
%% ===== END double-blind block =====

\begin{abstract}
Speech-to-speech translation (S2ST) for low-resource tonal languages on edge devices faces a triple constraint: limited compute, semantic sensitivity to tone, and high phrase repetition in industrial deployments. Naive caching strategies fail: byte-hash caches achieve 0\% hit rate on real speech, while generic audio fingerprints --- optimized for music identification --- are tone-blind and produce semantically catastrophic false hits on Vietnamese tonal minimal pairs such as \emph{ma} (ghost) / \emph{m\'a} (mother) / \emph{m\d{a}} (tomb). We propose a 272-dimensional \emph{tonal-perceptual fingerprint} composed of a tone histogram, temporal-subsampled MFCC, MFCC first-difference, F0 statistics, and energy statistics, and a \emph{two-stream voting} cache architecture that gates each hit on simultaneous agreement between an independent 6-d tone stream ($\tau_A=0.85$) and a 260-d segmental stream ($\tau_B=0.93$). On a synthesized Vietnamese tonal minimal-pair corpus, two-stream voting achieves 100\% correct-miss rate on tonal variants (vs.\ uniform false-hit failure for naive concatenation), while preserving 1{,}441$\times$ end-to-end speedup over an FP32 cascade on cache-hit ($25$~ms vs.\ $36.5$~s). The mechanism generalizes to any tonal language by adjusting the tone-histogram dimension to the target tone inventory (Mandarin: 5; Thai: 5; Yoruba: 3). We report honest negative results --- weighted-concat fingerprint fix, coarse-to-fine cascade routing, naive speaker-personalized cache --- to ground the design choice, and characterize cache scalability through $N=1024$ entries, where lookup remains fingerprint-bound rather than comparison-bound. The full system runs in 861~MB of model footprint, fitting Qualcomm Snapdragon 8 Gen 3 deployment for the OneVoice AI Challenge 2026 industrial-translator scenario.
\end{abstract}

\begin{highlights}
\item Linguistic-typology-aware fingerprint design that bridges speech and retrieval communities.
\item Two-stream voting cache architecture providing 100\% tonal-minimal-pair protection.
\item 1{,}441$\times$ end-to-end speedup on cache-hit while preserving WER and tone-error rate.
\item Cross-language generalization recipe via tone-inventory adjustment of fingerprint dimension.
\item Reproducible honest-failure log of three abandoned design alternatives.
\end{highlights}

\begin{keywords}
on-device speech translation \sep low-resource languages \sep tonal languages \sep Vietnamese \sep audio fingerprinting \sep edge AI \sep quantization \sep utterance caching
\end{keywords}

\maketitle

%% =====================================================================
%% Section 1 - Introduction
%% =====================================================================
\section{Introduction}\label{sec:intro}

% Motivation: low-resource tonal-language S2ST on edge devices for industrial workers
% OneVoice AI Challenge 2026: workers repeat short directive phrases many times per shift
% Existing caches:
%   - Exact byte-hash -> 0% hit rate on real speech
%   - Generic perceptual hash (music-ID) -> tone-blind, false-hit on minimal pairs
% Our contribution: 272-d tonal-perceptual fingerprint that
%   - Bridges linguistic typology and edge cache design
%   - Discriminates phrases (>5pp cosine gap) while admitting intra-speaker variation
%   - Generalizes to any tonal language via tone-histogram dim adjustment
%   - Composes with INT8 quantization and standard streaming MT into a deployable pipeline
% Headline number: 1,441x speedup on cache hit (25 ms vs 36.5 s FP32 cold)
% Roadmap

\section{Related Work}\label{sec:related}

\subsection{Speech translation}\label{sec:related:s2st}
% SeamlessM4T, NLLB, cascade vs E2E. Most prior systems target high-resource languages.

\subsection{On-device inference}\label{sec:related:edge}
% faster-whisper, CTranslate2, GGML. Engineering substrate, not novelty.

\subsection{Audio fingerprinting}\label{sec:related:fingerprint}
% Shazam-style \citep{wang2003shazam}, neural fingerprints. Optimized for music ID, not phrase-level cache.

\subsection{Cache architectures for ML}\label{sec:related:cache}
% BERT inference cache, KV-cache compression. Token-level, not utterance-level.

\subsection{Tonal-language speech}\label{sec:related:tonal}
% PhoWhisper \citep{le2024phowhisper}, de la Fuente \& Jurafsky 2024 Mandarin tone layer-probe \citep{delafuente2024tone}.

\subsection{Quantization for speech models}\label{sec:related:quant}
% Xu et al.\ 2025 \citep{xu2025mixedprecision}, FADE \citep{fade2026}, TAQ \citep{taq2025}. Our delta: typology-aware allocation.

% Gap: no prior work treats utterance-level caching for tonal speech translation.

%% =====================================================================
%% Section 3 - Method
%% =====================================================================
\section{Method}\label{sec:method}

\subsection{System overview}\label{sec:method:overview}
% Pipeline diagram: tone-VAD -> T1 audio cache -> ASR (INT8 streaming) -> T2 text cache -> MT (INT8) -> TTS (INT8)

\subsection{Theoretical framework}\label{sec:method:theory}
% Fingerprint function f: A -> R^D mapping audio to L2-normalized vector
% Goal 1 (invariance): cos(f(a), f(a')) >= tau for same-phrase realizations
% Goal 2 (discrimination): cos(f(a), f(a'')) < tau for different phrases
% Threshold tau as cache hit gate

\subsection{Tonal-perceptual fingerprint design}\label{sec:method:fingerprint}
% Five components: A. Tone histogram (6d), B. Temporal-subsampled MFCC (130d),
% C. MFCC first-difference (130d), D. F0 statistics (4d), E. Energy stats (2d). Total 272d.

\subsection{Text-level translation memory}\label{sec:method:textcache}
% 256-d diacritic-robust hashed-ngram fingerprint on ASR transcript. tau=0.80.

\subsection{Two-stream voting architecture}\label{sec:method:twostream}
% PAPER ANCHOR (contribution_analysis.md sec 2.7). Tone stream (6d, tau_A=0.85)
% + segmental stream (260d, tau_B=0.93). Hit iff both agree on same stored entry.
% 100% correct-miss rate on tonal minimal pairs vs uniform failure of naive concat.

\subsection{Industrial corpus pre-warming}\label{sec:method:prewarm}
% 42 curated Vi industrial phrases x 3 target languages loaded at startup.

\subsection{Tonal-language VAD}\label{sec:method:tonevad}
% F0-stability gate as typology-aware pre-filter.

\subsection{INT8 substrate}\label{sec:method:int8}
% faster-whisper + CTranslate2 NLLB + dyn-INT8 VITS. 3.9x size reduction, 5.4x speedup.

\subsection{Streaming progressive output}\label{sec:method:streaming}
% Generator-driven cascade emits partial text as decoders progress.

\subsection{Persistent cache snapshots}\label{sec:method:persist}
% JSON serialization with base64 float16 fingerprints + zlib int16 PCM audio.

%% =====================================================================
%% Section 4 - Experiments
%% =====================================================================
\section{Experiments}\label{sec:experiments}

\subsection{Datasets}\label{sec:exp:data}
% VIVOS \citep{luong2016vivos}, FLEURS-vi \citep{conneau2022fleurs}, AISHELL-3, DEMAND+MUSAN noise.

\subsection{Fingerprint-component ablation}\label{sec:exp:ablation}
% Table 1: strip each of A,B,C,D,E. Show MFCC (B) dominant, delta (C) +0.06 gap,
% tone histogram (A) 0 measurable gain on phrase-level VIVOS but critical on minimal pairs.

\subsection{Tonal minimal-pair evaluation}\label{sec:exp:minpair}
% 6 Vi syllables x 6 tones = 36 clips via Vi MMS-TTS carrier-frame synthesis.
% Tone-only fingerprint provides +0.028 protect gap; naive concat fails.

\subsection{Two-stream voting evaluation}\label{sec:exp:twostream}
% 30 clips, 100% correct-miss, 0 false hits.

\subsection{End-to-end latency ablation}\label{sec:exp:latency}
% Table 2: FP32, INT8, INT8+tone-VAD, INT8+T1, INT8+T2, INT8+T1+T2, INT8 full.

\subsection{Cache precision-recall calibration}\label{sec:exp:pr}
% threshold sweep tau in [0.70, 0.99]. Single-vector best F1=0.683 at tau=0.70.
% Operating point tau=0.95: P=0.958, R=0.115 (precision-first design for safety-critical).

\subsection{Cache scalability}\label{sec:exp:scale}
% N in {32, 128, 512, 1024}. Lookup is fingerprint-bound, not comparison-bound.

\subsection{Cross-language generalization}\label{sec:exp:crossling}
% Apply with tone-histogram dim = 5 to AISHELL-3 Mandarin. Verify gap preserved.

\subsection{Quality preservation}\label{sec:exp:quality}
% WER, CER, BLEU, chrF, COMET-22, MOS (n>=20 native raters).

%% =====================================================================
%% Section 5 - Deployment
%% =====================================================================
\section{Deployment Notes}\label{sec:deployment}
% Memory: 861 MB models + <=10 MB cache. Fits Snapdragon 8 Gen 3.
% Latency budget: <1.5 s end-to-end after tone-aware quantization.
% Energy: passive cooling.

%% =====================================================================
%% Section 6 - Limitations
%% =====================================================================
\section{Limitations and Negative Results}\label{sec:limitations}

\subsection{Speaker-personalized cache failure}\label{sec:lim:speaker}
% 26-d MFCC mean+std voice fingerprint cannot discriminate Vi speakers (cos 0.93-0.98).

\subsection{Coarse-to-fine MT failure}\label{sec:lim:c2f}
% opus-mt-vi-en produces empty translations on Vi industrial phrases.

\subsection{Initial global-MFCC fingerprint failure}\label{sec:lim:globalmfcc}
% Inter-phrase cos 0.99. Justifies temporal-subsampling design choice.

\subsection{Tone-VAD F0 range mis-calibration}\label{sec:lim:tonevad}
% 1.8 oct cap rejected real Vi speech (2.7 oct natural range). Re-tuned to 4.0 oct.

\subsection{Weighted-concat fingerprint fix failure}\label{sec:lim:weighted}
% All alpha weights gave negative protect-gap. Motivates two-stream voting.

%% =====================================================================
%% Section 7 - Conclusion
%% =====================================================================
\section{Conclusion}\label{sec:conclusion}

\clearpage

%% Placeholder figure
\begin{figure}
  \centering
  %\includegraphics{../figures/pipeline.pdf}
  \caption{System pipeline: tone-VAD $\rightarrow$ T1 audio cache $\rightarrow$ ASR (INT8 streaming) $\rightarrow$ T2 text cache $\rightarrow$ MT (INT8) $\rightarrow$ TTS (INT8). Two-stream voting gates T1 hits on simultaneous tone-stream and segmental-stream agreement.}\label{fig:pipeline}
\end{figure}

%% Placeholder table
\begin{table}
\caption{Fingerprint-component ablation on VIVOS (50 intra-phrase + 50 inter-phrase pairs). Higher gap is better.}\label{tab:fp_ablation}
\begin{tabular*}{\tblwidth}{@{}LLLLL@{}}
\toprule
Variant & Dims & Intra cos & Inter cos & Gap \\
\midrule
full (proposed) & 272 & 0.904 & 0.757 & 0.148 \\
no\_tone & 266 & 0.904 & 0.757 & 0.148 \\
no\_mfcc & 12 & 0.994 & 0.996 & -0.002 \\
no\_delta & 142 & 0.952 & 0.860 & 0.092 \\
mfcc\_only & 130 & 0.952 & 0.860 & 0.092 \\
tone\_only & 6 & 0.875 & 0.903 & -0.027 \\
\bottomrule
\end{tabular*}
\end{table}

%% Appendix
\appendix
\section{Pipeline and cache architecture diagram}\label{app:pipeline}
\section{Fingerprint component pseudocode}\label{app:pseudocode}
\section{Industrial phrase corpus}\label{app:phrases}
\section{Reproducibility checklist}\label{app:repro}

\printcredits

%% Loading bibliography style file
\bibliographystyle{cas-model2-names}

%% Loading bibliography database (placeholder entries in references.bib)
\bibliography{references}

\end{document}
